\documentclass[11pt]{article}
\usepackage[margin=1in]{geometry}
\usepackage{enumitem}
% =============================================================================
% Preambles/header.tex — shared LaTeX/Beamer preamble for this project
%
% Use from a Beamer lecture by prepending:
%     \input{header}
% and compiling with TEXINPUTS=../Preambles:$TEXINPUTS (the /compile-latex
% skill does this automatically).
%
% PALETTE CONTRACT. Color names defined here MUST match the names in
% Quarto/theme-template.scss so Beamer and Quarto renderings use the same
% palette. The scripts/check-palette-sync.sh script enforces this.
%
% To customize: edit the HEX values below AND the matching $variable in
% Quarto/theme-template.scss, then run scripts/check-palette-sync.sh.
% =============================================================================

% -----------------------------------------------------------------------------
% Engine detection + encoding
%
% Our /compile-latex skill uses XeLaTeX, where inputenc is unnecessary and
% can produce encoding warnings. Only load inputenc under pdfTeX.
% -----------------------------------------------------------------------------
\usepackage{iftex}
\ifPDFTeX
  \usepackage[utf8]{inputenc}
\fi

\usepackage{amsmath, amssymb, amsthm}
\usepackage{booktabs}
\usepackage{graphicx}

% -----------------------------------------------------------------------------
% PALETTE — mirrors Quarto/theme-template.scss variable names.
% Load xcolor and define colors BEFORE anything that references them
% (notably hyperref, hypersetup, and the Beamer color assignments).
% -----------------------------------------------------------------------------
\usepackage{xcolor}

\definecolor{primary-blue}{HTML}{012169}      % headings, accents
\definecolor{primary-gold}{HTML}{B9975B}      % emphasis, borders
\definecolor{highlight-yellow}{HTML}{F2A900}  % markers, alerts
\definecolor{light-bg}{HTML}{E8EDF5}          % title / section backgrounds
\definecolor{jet}{HTML}{1A1A1A}               % body text

% Semantic colors (used in diagrams and annotations)
\definecolor{positive}{HTML}{15803D}          % good / observed / identified
\definecolor{negative}{HTML}{B91C1C}          % bad / problematic / bias
\definecolor{neutral}{HTML}{525252}           % reference / context

% Highlight accents (match .hi-* classes in the SCSS)
\definecolor{hi-slate}{HTML}{314F4F}
\definecolor{hi-green}{HTML}{2E8B57}
\definecolor{hi-red}{HTML}{C41E3A}

% Semantic aliases so skill templates and snippets can write
% \textcolor{accent}{...} without hard-coding "primary-blue".
\colorlet{accent}{primary-blue}
\colorlet{accent2}{primary-gold}

% -----------------------------------------------------------------------------
% Hyperref — loaded AFTER xcolor + palette so link colors resolve.
% -----------------------------------------------------------------------------
\usepackage{hyperref}
\hypersetup{colorlinks=true, linkcolor=primary-blue, urlcolor=primary-blue,
            citecolor=primary-blue, pdfborder={0 0 0}}

% -----------------------------------------------------------------------------
% Course code — set once per deck (after \input{header}, before \begin{document})
% via \coursecode{CS401}. Shown in the footer on every slide so a deck's course
% is identifiable without opening the title page. Empty by default.
% -----------------------------------------------------------------------------
\makeatletter
\newcommand{\coursecode}[1]{\gdef\@coursecodetext{#1}}
\gdef\@coursecodetext{}
\makeatother

% -----------------------------------------------------------------------------
% Beamer theme assignments (only apply under Beamer).
%
% We detect Beamer via \@ifclassloaded{beamer}, which is the documented,
% non-internal way. This must be wrapped in \makeatletter/\makeatother
% because \@ifclassloaded contains an @-catcode character.
% -----------------------------------------------------------------------------
\makeatletter
\@ifclassloaded{beamer}{%
  \usefonttheme{professionalfonts}
  \setbeamercolor{structure}{fg=primary-blue}
  \setbeamercolor{frametitle}{fg=primary-blue}
  \setbeamercolor{title}{fg=primary-blue}
  \setbeamercolor{subtitle}{fg=primary-gold}
  \setbeamercolor{author}{fg=primary-blue}
  \setbeamercolor{institute}{fg=neutral}
  \setbeamercolor{date}{fg=primary-gold}
  \setbeamercolor{itemize item}{fg=primary-blue}
  \setbeamercolor{itemize subitem}{fg=primary-gold}
  \setbeamercolor{alerted text}{fg=primary-gold}
  \setbeamercolor{block title}{fg=white, bg=primary-blue}
  \setbeamercolor{block body}{bg=light-bg}
  % Semantic block variants: block=definition/concept (blue), exampleblock=
  % worked example (green/positive), alertblock=key takeaway or caution (gold).
  % Keep to <=2 colored boxes per slide across ALL three variants (INV-7).
  \setbeamercolor{block title example}{fg=white, bg=positive}
  \setbeamercolor{block body example}{bg=light-bg}
  \setbeamercolor{block title alerted}{fg=white, bg=primary-gold}
  \setbeamercolor{block body alerted}{bg=light-bg}

  % Minimal footer: course code (left) + page X / N (right), no navigation clutter
  \setbeamertemplate{navigation symbols}{}
  \setbeamertemplate{footline}{%
    \color{neutral}\footnotesize\hspace{0.7em}\@coursecodetext\hfill
    \insertframenumber/\inserttotalframenumber\hspace{0.7em}\vspace{0.3em}}
}{}
\makeatother

% -----------------------------------------------------------------------------
% TikZ setup — libraries the snippet gallery uses
% -----------------------------------------------------------------------------
\usepackage{tikz}
\usetikzlibrary{arrows.meta, positioning, calc, decorations.pathreplacing,
                fit, shapes.geometric, backgrounds}

% Shared TikZ styles — snippets and hand-written diagrams can pick these up
% via `\begin{tikzpicture}[every node/.style={...}]` or by naming specific
% styles (e.g., `\node[dag-node] (X) at (0,0) {X};`).
\tikzset{
  % Node styles for causal DAGs and similar
  dag-node/.style = {
    draw=primary-blue, fill=light-bg, rounded corners=2pt,
    minimum width=1.2cm, minimum height=0.9cm, align=center, font=\small
  },
  decision-node/.style = {
    draw=primary-gold, fill=white, diamond, aspect=1.8,
    minimum width=1.8cm, minimum height=1.2cm, align=center, font=\small,
    inner sep=1pt
  },
  % Edge styles — observed vs counterfactual
  observed-edge/.style = {-{Latex[length=2.5mm]}, thick, draw=jet},
  counterfactual-edge/.style = {-{Latex[length=2.5mm]}, thick, draw=neutral, dashed},
  confound-edge/.style = {-{Latex[length=2.5mm]}, thick, draw=negative, dashed},
  % Data-point styles for plots
  observed-dot/.style = {fill=primary-blue, draw=primary-blue, circle, inner sep=1.2pt},
  counterfactual-dot/.style = {fill=white, draw=neutral, circle, inner sep=1.2pt}
}

% -----------------------------------------------------------------------------
% Convenience macros
% -----------------------------------------------------------------------------
\newcommand{\muted}[1]{\textcolor{neutral}{#1}}
\newcommand{\key}[1]{\textcolor{primary-gold}{\textbf{#1}}}
\newcommand{\good}[1]{\textcolor{positive}{#1}}
\newcommand{\bad}[1]{\textcolor{negative}{#1}}

% Standout frame macro: full-bleed dark background for section transitions.
% Beamer-only (uses \begin{frame}/\setbeamercolor/\usebeamerfont) -- do not
% call from an article-class file (e.g. Notes/<CODE>/*-notes.tex). \muted,
% \key, \good, \bad, and \coursecode above are plain xcolor/gdef and are
% safe to use from either class; verified when Notes/article-class support
% was added.
% Expand: \transitionslide{Your transition title}
\newcommand{\transitionslide}[1]{%
  \begingroup
  \setbeamercolor{background canvas}{bg=primary-blue}
  \begin{frame}[plain]
    \centering
    \vfill
    {\usebeamerfont{title}\color{white}\Large #1\par}
    \vfill
  \end{frame}
  \endgroup
}

% Section divider frame (Beamer-only), an alternative to \transitionslide with a
% darker two-line style: near-black (jet) background, a small white label line
% above a larger gold title, and a thin gold rule along the bottom edge.
% Expand: \sectiondivider{Section 2}{Pointers}
\newcommand{\sectiondivider}[2]{%
  \begingroup
  \setbeamercolor{background canvas}{bg=jet}
  \begin{frame}[plain]
    \vspace*{3.0cm}
    {\color{white}\ttfamily\large #1\par}
    \vspace{0.5cm}
    {\color{primary-gold}\LARGE #2\par}
    \begin{tikzpicture}[remember picture, overlay]
      \draw[primary-gold, line width=2.5pt]
        ($(current page.south west)+(0,0.1)$) -- ($(current page.south east)+(0,0.1)$);
    \end{tikzpicture}
  \end{frame}
  \endgroup
}

% -----------------------------------------------------------------------------
% Channel watermark — "Concepts That Click" (YouTube).
%
% Article-class files (CompetitiveExam Q/A sets, Notes, Assignments) get a
% light diagonal draft watermark on every page plus a centered footer naming
% the channel. Beamer decks get a subtle TikZ background watermark instead
% (the footline already carries the course code). Centralized here so every
% MCQ PDF is branded without per-file edits.
% -----------------------------------------------------------------------------
\makeatletter
\@ifclassloaded{article}{%
  \usepackage{draftwatermark}
  \SetWatermarkText{Concepts That Click}
  \SetWatermarkScale{1}
  \SetWatermarkFontSize{2cm}
  \SetWatermarkLightness{0.86}
  \SetWatermarkAngle{45}
  \usepackage{fancyhdr}
  \pagestyle{fancy}
  \fancyhf{}
  \fancyfoot[C]{\footnotesize\textcolor{neutral}{Concepts That Click~|~YouTube: \href{https://www.youtube.com/@concepts.thatclick}{@concepts.thatclick}}}
  \renewcommand{\headrulewidth}{0pt}
  \renewcommand{\footrulewidth}{0pt}
  \fancypagestyle{plain}{%
    \fancyhf{}%
    \fancyfoot[C]{\footnotesize\textcolor{neutral}{Concepts That Click~|~YouTube: \href{https://www.youtube.com/@concepts.thatclick}{@concepts.thatclick}}}%
    \renewcommand{\headrulewidth}{0pt}%
    \renewcommand{\footrulewidth}{0pt}%
  }
}{}
\@ifclassloaded{beamer}{%
  \setbeamertemplate{background}{%
    \begin{tikzpicture}[remember picture, overlay]
      \node[rotate=45, opacity=0.055, align=center]
        at (current page.center)
        {\Huge\textcolor{neutral}{Concepts That Click}};
    \end{tikzpicture}%
  }
}{}
\makeatother 
\coursecode{JKSSB}
\renewcommand{\thesection}{22.\arabic{section}}
\newtheorem{example}{Example}[section]
\newenvironment{definitionbox}[1]{%
  \par\noindent\textbf{Definition (#1).}\ \itshape
}{\par\vspace{0.3em}}
\title{MS Word Fundamentals --- Detailed Lecture Notes}
\author{JKSSB Computer Awareness}
\date{\today}

\begin{document}
\maketitle

\section{What MS Word is}
\textbf{MS Word} (Microsoft Word) is a \textbf{word processor}: a program
used to create, edit, format, store, and print text documents. It is one of
the main applications in the \textbf{Microsoft Office} suite. A document
written in Word is saved as a file, and the modern default file format uses
the \textbf{.docx} extension. Word shows the document in a rectangular
window that contains the document area together with several command areas
around it.

\begin{definitionbox}{Word processor}
An application that lets a user type text, change its appearance
(formatting), arrange it on pages, and print it. MS Word is the common
example.
\end{definitionbox}

A word processor separates the \emph{text} from its \emph{appearance}. The
same words can be displayed in many different fonts, sizes, colours, and
layouts, because the formatting is stored alongside the text and can be
changed at any time. This separation is the central idea of the whole
lesson.

\begin{center}
\begin{tabular}{p{0.28\textwidth}p{0.60\textwidth}}
\toprule
\textbf{Part of the window} & \textbf{What it is} \\
\midrule
Document area & The white page where the text is typed and shown. \\
Ribbon & The strip of command tabs across the top of the window. \\
Quick Access Toolbar & The small row of always-visible buttons at the top-left. \\
Status bar & The strip at the bottom showing page, word count, and view buttons. \\
\bottomrule
\end{tabular}
\end{center}

\section{The Ribbon and its tabs}
The \textbf{Ribbon} is the broad band of commands across the top of the Word
window. It replaced the older style of drop-down menus and toolbars. The
Ribbon is organised in three levels:

\begin{enumerate}
  \item \textbf{Tabs} across the top: \textbf{Home, Insert, Design, Layout,
    References, Mailings, Review,} and \textbf{View}.
  \item \textbf{Groups} inside each tab. For example, the Home tab contains
    the \textbf{Font}, \textbf{Paragraph}, and \textbf{Styles} groups.
  \item \textbf{Commands} inside each group, such as the Bold button or the
    alignment buttons.
\end{enumerate}

Some tabs are \textbf{contextual}. They are not shown all the time; they
appear only when the right object is selected. When a table is selected, the
\textbf{Table Tools} tabs appear; when a picture is selected, the
\textbf{Picture Format} tab appears. The \textbf{File} tab is special: it
does not show groups like the other tabs but opens \textbf{Backstage view},
which holds document-level actions such as New, Open, Save, Save As, Print,
and Close.

\begin{definitionbox}{Ribbon}
The command area of the Word window, organised into tabs, groups, and
buttons, which replaced the older menu-and-toolbar layout.
\end{definitionbox}

For exam purposes, the key point is that the \textbf{Home} tab carries the
most-used formatting commands: the \emph{Font} tools that change how
characters look, and the \emph{Paragraph} tools that change alignment,
indentation, spacing, bullets, and numbering.

\section{The Quick Access Toolbar}
The \textbf{Quick Access Toolbar (QAT)} is the small row of buttons in the
top-left corner of the window, just above the Ribbon. It is designed to hold
the few commands a user needs most often, so they stay visible on every tab.
By default it shows \textbf{Save}, \textbf{Undo}, and \textbf{Redo}. A user
can \textbf{customise} it by adding or removing commands, and it can be moved
to sit below the Ribbon if preferred.

\begin{definitionbox}{Quick Access Toolbar}
A small, customisable toolbar at the top-left of the Word window that keeps
frequently used commands (by default Save, Undo, Redo) visible on every tab.
\end{definitionbox}

\section{Creating, opening and saving documents}
The \textbf{File} tab opens \textbf{Backstage view}, the part of Word where
you work with the document as a whole rather than its contents. From here a
user can create a new document, open an existing one, save, print, or close.
The common keyboard shortcuts are worth memorising because direct-recall
items ask for them.

\begin{center}
\begin{tabular}{p{0.24\textwidth}p{0.64\textwidth}}
\toprule
\textbf{Shortcut} & \textbf{Action} \\
\midrule
Ctrl+N & Create a new blank document. \\
Ctrl+O & Open an existing document. \\
Ctrl+S & Save the current document. \\
Ctrl+P & Open the Print screen. \\
\bottomrule
\end{tabular}
\end{center}

\textbf{Save} writes the current document back to its file. \textbf{Save As}
writes a \emph{copy} under a new name, in a new location, or in a different
file format, leaving the original as it was. The default Word format is
\textbf{.docx}; documents made with much older versions use the \textbf{.doc}
extension.

\begin{example}
A student types a letter and presses Ctrl+S, saving it as
\texttt{letter.docx}. To keep the original and also produce a version for
printing, the student uses Save As and saves a second file with a new name.
The first file is unchanged.
\end{example}

\section{Entering and editing text}
Text is typed at the \textbf{insertion point}, which is shown by a blinking
cursor. \textbf{Word wrap} moves text automatically to the next line when
the current line is full; pressing \textbf{Enter} starts a new
\textbf{paragraph}. The \textbf{Backspace} key deletes the character to the
left of the cursor, while the \textbf{Delete} key deletes the character to
the right. \textbf{Ctrl+Z} undoes the last action and \textbf{Ctrl+Y} redoes
it.

Text must be selected before it can be reformatted. A \textbf{word} is
selected by double-clicking it, a \textbf{paragraph} by triple-clicking, and
the \textbf{whole document} with \textbf{Ctrl+A}.

\section{Font formatting}
The \textbf{Font} group on the Home tab changes how individual characters
look. It sets the \textbf{font face} (the typeface, such as Arial or Times
New Roman) and the \textbf{font size}. It provides the style buttons
\textbf{Bold (Ctrl+B)}, \textbf{Italic (Ctrl+I)}, and \textbf{Underline
(Ctrl+U)}, and it sets the \textbf{font colour} and \textbf{highlight}. It
also applies effects such as subscript and superscript. Deeper options are
available in the Font dialog box.

\begin{definitionbox}{Font}
The typeface and its settings (face, size, style, colour) that determine how
characters are displayed. Font formatting changes the appearance of
characters, not their arrangement on the page.
\end{definitionbox}

\section{Paragraph alignment}
\textbf{Alignment} decides where each line of a paragraph sits between the
left and right margins.

\begin{center}
\begin{tabular}{p{0.22\textwidth}p{0.20\textwidth}p{0.46\textwidth}}
\toprule
\textbf{Alignment} & \textbf{Shortcut} & \textbf{Effect} \\
\midrule
Left & Ctrl+L & Text lines up evenly on the left; ragged on the right. This is the default. \\
Centre & Ctrl+E & Each line is placed in the middle of the page. \\
Right & Ctrl+R & Text lines up evenly on the right. \\
Justify & Ctrl+J & Text is flush with both left and right margins. \\
\bottomrule
\end{tabular}
\end{center}

\textbf{Justify} stretches the spaces between words so that both edges of
the paragraph are even, which is why books and newspapers use it.
\textbf{Centre} does not stretch anything; it simply moves each line to the
middle.

\section{Indentation}
\textbf{Indentation} pushes paragraph lines inward, away from the margin.
The \textbf{left indent} and \textbf{right indent} move the whole paragraph
inward from the respective side. A \textbf{first-line indent} pushes only
the first line inward, as in a typical book paragraph. A \textbf{hanging
indent} pushes every line \emph{except} the first inward, which is the
standard layout for a references or bibliography list.

\begin{definitionbox}{Indentation}
The inward offset of paragraph lines from the page margins. First-line
indent affects only the first line; hanging indent affects every line except
the first.
\end{definitionbox}

Indentation must not be confused with \textbf{spacing}. Indentation shifts
lines \emph{horizontally}; spacing changes the \emph{vertical} gaps.

\section{Line and paragraph spacing}
\textbf{Line spacing} sets the vertical gap between lines \emph{inside} a
paragraph; common settings are single, 1.5, and double. \textbf{Paragraph
spacing} sets the gap \emph{before} and \emph{after} a paragraph. Both are
controlled from the Paragraph group on the Home tab, or from the Paragraph
dialog box. Using spacing is better than pressing Enter repeatedly to make
gaps, because spacing stays correct when text is edited.

\section{Bullets and numbering}
\textbf{Bullets} mark list items with a symbol and are used when the order of
the items does not matter. \textbf{Numbering} marks items with numbers or
letters and is used when the order or sequence matters. \textbf{Multilevel
lists} combine levels (for example 1, then 1.1, then 1.1.1) to build nested
outlines. All three are applied from the \textbf{Paragraph} group on the
Home tab.

\begin{example}
A shopping list uses \textbf{bullets} because the order is irrelevant. A set
of instructions to install software uses \textbf{numbering} because the
steps must be followed in order.
\end{example}

\section{Tables}
A \textbf{table} arranges information in \textbf{rows} and \textbf{columns};
each intersection of a row and a column is a \textbf{cell}. A table is
inserted from \textbf{Insert $\rightarrow$ Table}. Cells can be
\textbf{merged} into a single larger cell or \textbf{split} back into
separate cells, and rows and columns can be inserted, deleted, and resized.
When the table is selected, the \textbf{Table Tools} contextual tabs appear.

\begin{definitionbox}{Table}
A grid of rows and columns used to arrange text in cells, inserted in Word
from Insert $\rightarrow$ Table.
\end{definitionbox}

\section{Images and objects}
An image is placed in the document from \textbf{Insert $\rightarrow$
Pictures}. \textbf{Handles} around the image let the user resize and rotate
it. \textbf{Text wrapping} decides how the surrounding text flows around the
image; options include In Line with Text, Square, Tight, and Behind Text.
When the image is selected, the \textbf{Picture Format} contextual tab
appears. Charts, shapes, and \textbf{SmartArt} are inserted from the same
Insert tab.

\begin{definitionbox}{Text wrapping}
The setting that controls how text flows around a picture or object. ``In
Line with Text'' treats the image like a character; the other styles let
text flow around the image.
\end{definitionbox}

\section{Headers and footers}
A \textbf{header} is the area at the \emph{top} of a page and a
\textbf{footer} is the area at the \emph{bottom}. Both repeat on every page
of a section, which makes them ideal for a document title (in the header)
and \textbf{page numbers} (in the footer). They are inserted from
\textbf{Insert $\rightarrow$ Header} and \textbf{Insert $\rightarrow$
Footer}.

\section{Page setup}
Page setup is found on the \textbf{Layout} tab. It controls the overall page
geometry: \textbf{margins} (the blank space between the text and the page
edges), \textbf{orientation} (\textbf{Portrait}, which is tall, or
\textbf{Landscape}, which is wide), and \textbf{page size} (for example A4
or Letter). The same tab sets \textbf{columns} and the page \textbf{breaks}
described next.

\section{Page break vs section break}
A \textbf{Page Break} (\textbf{Ctrl+Enter}) forces the following content to
start on a \emph{new page}, but the page formatting stays the same. A
\textbf{Section Break} starts a new \textbf{section} of the document, which
can carry its own page setup. The most important consequence for the exam is
that only a \textbf{Section Break} allows a part of the document to have
\textbf{different headers and footers}. Both are inserted from
\textbf{Layout $\rightarrow$ Breaks}.

\begin{example}
An item claims, ``A Page Break can be used to give the first page a
different header from the rest of the document.'' This is wrong. A Page
Break only moves content to a new page; to change the header or footer part
way through the document, a \textbf{Section Break} is required. The wrong
option has confused the two breaks (TRAP-13).
\end{example}

\section{Printing a document}
\textbf{Ctrl+P} opens the \textbf{Print} screen in Backstage view. The
\textbf{print preview} shows exactly how the pages will appear on paper. On
this screen the user selects the \textbf{printer}, the number of
\textbf{copies}, and which \textbf{pages} to print, then clicks \textbf{Print}
to send the job to the printer. Checking the preview first catches wrong
margins, stray blank pages, and cut-off text before paper is used.

\section{Exam retrieval and common confusions}
Five distinctions prevent most errors on this topic.
\begin{enumerate}
  \item \textbf{Save} overwrites the current file; \textbf{Save As} stores a
    copy under a new name, location, or format.
  \item \textbf{Ctrl+S} saves; \textbf{Ctrl+P} prints; \textbf{Ctrl+Enter}
    inserts a page break. Do not swap the shortcuts.
  \item \textbf{Bullets} are for unordered lists; \textbf{numbering} is for
    ordered lists.
  \item \textbf{Indentation} shifts lines horizontally; \textbf{spacing}
    changes vertical gaps between lines and paragraphs.
  \item A \textbf{Page Break} does \emph{not} change headers and footers;
    only a \textbf{Section Break} does (TRAP-13).
\end{enumerate}

\begin{example}
An item asks which shortcut opens the Print screen. Trace the File-level
actions: New is Ctrl+N, Open is Ctrl+O, Save is Ctrl+S, and Print is
Ctrl+P. The answer is \textbf{Ctrl+P}. If the stem instead asks for Undo, the
answer is Ctrl+Z, and if it asks for a page break, the answer is Ctrl+Enter.
\end{example}

\subsection*{Exercises}
\begin{enumerate}[label=\arabic*.]
  \item Name the main parts of the MS Word window.
  \item What is the Quick Access Toolbar, and what does it show by default?
  \item Distinguish Save from Save As.
  \item List the four paragraph alignments with their shortcuts.
  \item Explain the difference between a Page Break and a Section Break.
\end{enumerate}

\subsection*{Solutions}
\begingroup\small
\begin{enumerate}[label=\arabic*.]
  \item The document area, the Ribbon (with its tabs and groups), the Quick
    Access Toolbar, and the status bar.
  \item The Quick Access Toolbar is the small customisable row of buttons at
    the top-left of the window. By default it shows Save, Undo, and Redo,
    and it stays visible on every tab.
  \item Save writes the current document back to its existing file. Save As
    writes a copy under a new name, location, or format and leaves the
    original file unchanged.
  \item Left (Ctrl+L), Centre (Ctrl+E), Right (Ctrl+R), and Justify
    (Ctrl+J).
  \item A Page Break (Ctrl+Enter) starts the following content on a new page
    while keeping the same formatting. A Section Break starts a new section
    that can have its own page setup --- including different headers and
    footers. Only a Section Break can change headers and footers part way
    through a document.
\end{enumerate}
\endgroup
\end{document}
